\scanpage{3}\mspage{2}

Notons que $M_g$ et $M_{g'}$ définissent le \={m} élément de $\underline{\mathrm{NS}}_{P^\delta/S}(S)$ en vertu de b), or on a un isom canonique

\[
\underline{\mathrm{NS}}_{P^\delta/S}\simeq\underline{\mathrm{NS}}_{P^{\unsure{\circ|0}}/S}.
\]

Donc on trouve une section

\[
\varphi(\delta)\in\underline{\mathrm{NS}}_{P^\circ/S}(S)
\]

indépendante du choix de la section $\deleted{\unsure{\delta|g}}g$ de $X/S$. Par descente, elle est définie \={m} sans supposer qu'il existe une section. Comme la formation de $\varphi(\delta)$ est compatible (par construction) à l'extension de la base, on trouve en fait un homomorphisme de schémas (pas nécessairement additif)

% slanted
\margin{\uline{détailler}}

\[
\boxed{\varphi:\underline{\mathrm{NS}}_{X/S}\longrightarrow\underline{\mathrm{NS}}_{\deleted{P}B/S}}\qquad(B=P^\circ=\underline{\mathrm{Pic}}^\circ_{X/S}).
\]

d) Revenons à $\delta$\deleted{\illegible{1}}. Notons que $\underline{\mathrm{Pic}}_{P^\delta/S}$ est déduit de $\underline{\mathrm{Pic}}_{\deleted{P}B/S}$ en tordant par \deleted{l'action} le torseur $P^\delta$ sous $P^\circ$, via l'action naturelle de $\deleted{P}B$ sur $\underline{\mathrm{Pic}}_{\deleted{P}B/S}$ déduite de l'action de $\deleted{P}B$ sur lui-= par translation. \deleted{\unsure{Si $\lambda$|Si le}}

Cette action est triviale sur $\underline{\mathrm{Pic}}^\circ_{\deleted{P}B/S}$ et $\underline{\mathrm{NS}}_{B/S}$ [d'où en particulier un isom

\[
\underline{\mathrm{Pic}}^\circ_{P^\delta/S}\simeq\underline{\mathrm{Pic}}^\circ_{B/S},\qquad
\underline{\mathrm{NS}}_{P^\delta/S}\simeq\underline{\mathrm{NS}}_{B/S}\quad].
\]
